\section{Related work}
\paragraph{Language and visual evidence in Gaussian scenes.}
VALA addresses semantic assignment through visibility-aware lifting and robust
multi-view aggregation~\cite{vala}. Occam's LGS~\cite{occam} supplies a strong Gaussian
language baseline; both its published numbers and our retained reproduction
appear in the comparisons. LUDVIG lifts visual features or masks and uses graph
diffusion to improve scene queries~\cite{ludvig}. These methods motivate preserving
effective providers. Our variable is the decision among their complete answers,
which is separate from improving feature lifting or spatial smoothing.

EviSplat~\cite{evisplat} preserves individual observation features within
instances and combines query-selected instance evidence with local Gaussian
support. Query-dependent evidence use is therefore an established direction.
Our specific question concerns the utility of replacing an already complete
answer, assessed through intact-candidate and retention controls rather than
through representation reconstruction.

\paragraph{Verification and selective decisions.}
Our estimator predicts actual relative answer utility from observable evidence.
Its scalar candidate potential makes pairwise preferences cycle-consistent;
the global margin controls when a replacement is accepted. Neither property
guarantees accurate transfer. The decisive comparisons hold providers and
candidates fixed while removing query information or retention. Published
method comparisons supply context, while those matched controls attribute any
gain to the response decision. Input access, geometry and metric domains are
listed explicitly because they can differ between published methods and our
retained providers.
